
\documentclass[final, 12pt, aspectratio=169, xcolor={dvipsnames}]{beamer}
\usepackage{lmodern}
\usepackage[utf8]{inputenc}
\usepackage[T1]{fontenc}
\usepackage{graphicx}
\usepackage{xcolor}
\usepackage{textpos} % package for the positioning
\usepackage{tcolorbox}
\usepackage{tabularx}
\usepackage{hanging}
\usepackage{animate}
\usepackage{caption}
\usepackage{booktabs}
\usepackage[
    backend=biber,
    style=chicago-authordate-trad,
    sortlocale=de_DE,
    maxbibnames=20,
    natbib=true,
    url=false, 
    doi=false,
    eprint=false
]{biblatex}
\addbibresource{eureka.bib}

% make bibliography entries smaller
\renewcommand\bibfont{\scriptsize}

% Now get rid of all the colours
\setbeamercolor*{bibliography entry title}{fg=black}
\setbeamercolor*{bibliography entry author}{fg=black}
\setbeamercolor*{bibliography entry location}{fg=black}
\setbeamercolor*{bibliography entry note}{fg=black}

\title[PEARL]{Internal migration in the PBL/CBS Regional Population Projections}
\subtitle[PEARL]{}
\author[T. Husby]{Trond Grytli Husby, Andries de Jong\ }
\institute[PBL]{
  Netherlands Environmental Assesment Agency (PBL) \\[5ex]
  \texttt{trond.husby@pbl.nl}
}
\date[\today]{Amsterdam \today}

% set path to figures
\newcommand*{\figs}{../figs}%

% position the logo
\addtobeamertemplate{background}{}{%
  %\begin{textblock*}{100mm}(0.95\textwidth, 7cm)
  %\begin{flushright}
    \includegraphics[width=\paperwidth, height=\paperheight]{\figs/pbl_background.pdf}
  %\end{flushright}
    %\end{textblock*}}
}

% colour scheme and settings
\setbeamercolor{title}{bg=white,fg=blue!35!black}
\setbeamercolor{frametitle}{bg=,fg=PineGreen}
\setbeamercolor{enumerate item}{fg=PineGreen}
\setbeamercolor{itemize item}{fg=PineGreen}
\setbeamertemplate{itemize item}[circle]
\setbeamercolor{itemize subitem}{fg=PineGreen}
\setbeamertemplate{itemize subitem}[triangle]
\setbeamerfont{frametitle}{size=\normalsize}
\addtobeamertemplate{frametitle}{\vspace*{1cm}}{\vspace*{0.0cm}}
\setbeamertemplate{footnote}{\hangpara{2em}{1}\makebox[2em][l]{\insertfootnotemark}\footnotesize\insertfootnotetext\par}

% miscellaneous
\newcommand{\semitransp}[2][35]{\color{fg!#1}#2}
\newcommand{\source}[1]{\caption*{\tiny Source: {#1}} }


\begin{document}

\beamertemplatenavigationsymbolsempty

%--- the titlepage frame -------------------------%
{
  \setbeamertemplate{footline}{}

  \begin{frame}
    \titlepage
  \end{frame}
}

%--- frame 1 -------------------------%
\begin{frame}{The regional population projections by PBL and CBS}
  \begin{itemize}
    \item  Projections of population, households and demographic events in Dutch municipalities until 2050
  \item Carried out every three years: previous edition was in 2019
  \item Regional projections are made consistent with the national projections (CBS) 
  \item The projections are made with a combination of a cohort-component model (Projecting population Events at Regional Level), multistate life table model and internal migration module
    \item The projection model is called PEARL: projecting population events at a regional level
  \end{itemize}
\end{frame}

%--- frame 2 --------------------------------%
\begin{frame}{Structure of PEARL}
  \begin{figure}
      \includegraphics[trim={0 0 0 0.5cm},clip, scale = 0.4]{\figs/pearl_structure.png}
    \end{figure}
\end{frame}


%--- frame 3 --------------------------------%
\begin{frame}{Internal migration: arrivals}
  %\small
  \begin{itemize}
  \item Long-distance migration determined on the basis of observed long-distance patterns
  \item More than half of migration flows are short-distance relocations (<35 km). In PEARL this is represented with a \textit{constrained gravity model}
  \item The (initial) distribution of flows from municipality $i$ to $j$ is modelled as
    \begin{equation}
      M_{i,j} = O_{i}A_{i}\prod\limits^{k}X_{k,i,j}^{\hat{\beta}_{k,i}}D_{i,j}^{\hat{\gamma}_{i}},
      \quad O_{i} = \sum\limits_{j} M_{i,j},
      \quad A_{i} = \frac{1}{\sum\limits_{j}\prod\limits^{k}X_{k,i,j}^{\hat{\beta}_{k,i}}D_{i,j}^{\hat{\gamma}_{i}}}
    \end{equation}
  \item $O_{i}$ is determined in a previous step in PEARL
    \item $\hat{\beta}_{k,i}$ and $\hat{\gamma}_{i}$ are estimated origin-specific parameters
  \end{itemize}
\end{frame}

%--- frame 4 --------------------------------%
\begin{frame}{Research Question}
  %\small
  \begin{itemize}
  \item Earlier work focused on local modelling of origin-constrained gravity with \textit{geographically weighted regression (GWR)}
  \item Empirical strategy: take logs of both sides, rewrite, and estimate with OLS 
  \item \textbf{Which specification minimises \textit{out-of-sample} prediction errors? We compare results from a count model with OLS}
  \item Constraints are usually explicitly modelled, but with certain count models they can be captured by fixed effects. \textbf{Does this also apply if they are estimated using GWR?}
  \end{itemize}
\end{frame}


\begin{frame}{Origin-constrained gravity using Poisson regression}
  %\small
    $$ E(M_{i,j}) = \lambda_{i,j}  = \exp(\beta_{0, i} + \sum\limits_{k=1}^{K} \beta_{k} \log{X}_{k,i,j} + \gamma \log{D}_{i,j}) $$
  \begin{itemize}
    \item Maximum likelihood estimation with Iteratively Reweighted Least Squares (IRLS)
    \item Fixed effects $\beta_{0, i}$ ensure that the origin constraints hold: $\sum\limits_{j} \lambda_{i,j} = \sum\limits_{j} M_{i,j}$
    \end{itemize}
\end{frame}


\begin{frame}{Geographically weighted Poisson regression: origin-specific and origin-constrained}
  $$\lambda_{i,j} = \exp(\beta_{0,i} + \sum\limits_{k=1}^{K}\beta_{k}(\mathbf{u_{i}})\log{X_{k,i,j}} +  \gamma(\mathbf{u_{i}}) \log{D}_{i,j})$$
  \begin{itemize}
  \item Origins $i$ are represented by the population-weighted centroids of each municipality 
  \item The $k$th parameter for location $i$, $\beta_{k}(\mathbf{u}_{i})$ is a function of the coordinates $\mathbf{u}_{i}$
  \item $\beta_{0,i}$ is the (unweighted) fixed effects of $i$: origin constraints hold!
  \item Estimation with local IRLS \citep{nakaya2001local, nakaya2005geographically}, using a customised version of the \textit{GWmodel} package in \textit{R}: for estimation  we make use of sparsity and the network structure of the data
  \end{itemize}
\end{frame}
  
\begin{frame}{Geographically weighted Poisson regression}
  \begin{minipage}[t]{0.55\linewidth}%
    \footnotesize
    \begin{itemize}
    \item Weighted regression for each origin $i$ with distance-based weights according to a spatial kernel
    \item This creates a (weighted) neighbourhood around each origin
      \item Bandwidth is either distance or number of neighbouring points
    \item Model fitting completed with selection of kernel type and bandwith using cross validation. As a bonus, this allows us to simulate \textit{within-} and \textit{out-of-sample} prediction error, and to investigate the interplay between the two
      \end{itemize}
  \end{minipage}
  \begin{minipage}[t]{0.35\linewidth}%
    \vspace{-1cm}
    \begin{figure}
      \includegraphics[trim={2.8cm 0 0 0},clip, scale = 0.25]{\figs/kernels.png}
    \end{figure}
  \end{minipage}
\end{frame}

\begin{frame}{Data and variables}
  \begin{minipage}[t]{0.40\linewidth}%
    \footnotesize
    \begin{itemize}
      \item Short distance: $M_{i,j}$ where Euclidian distance between $i$ and $j$ is less than 35 km
      \item Network with 390 vertices (municipalities) and 14558 edges (bilateral origin destination flows)
      \item 2016 used as hold out sample. Evaluation of prediction accuracy by plugging estimated parameters into Equation 1
    \end{itemize}
  \end{minipage}
  \begin{minipage}[t]{0.48\linewidth}%
  \scriptsize
    \begin{table}{Explanatory variables}
      \begin{tabular}{ll}
        \hline
        Population  & $P_{j}$\\
        Road distance (km)    & $D_{i,j}$ \\
        Centrality   & $C_{i, j}$ = $\sum\limits_{l=1; l \ne i; l \ne j}^{L} \frac{P_{l}}{D_{j,l}}$\\
        Index of net new dwellings  & $H_{j} =  1 + \frac{\text{Net growth supply}_{j}}{\text{Stock}_{j}} $ \\
        \hline
        \end{tabular}
    \end{table}
    
    \vspace{-0.6cm}
    
    \begin{figure}
      \flushleft
    \includegraphics[trim={0 0 0 0},clip, scale = 0.3]{\figs/data_split.pdf}
  \end{figure}
  \end{minipage}
\end{frame}

\begin{frame}{Finding the optimal kernel with cross validation}
  \hspace{-0.5cm}
  \begin{minipage}[t]{0.68\linewidth}%
    \vspace{-0.3cm}
    \begin{figure}
      \includegraphics[trim={0 0 0 0},clip, scale = 0.5]{\figs/cv_oos_ws_comparison.png}
    \end{figure}
  \end{minipage}
  \begin{minipage}[t]{0.28\linewidth}%
    {\footnotesize
      %\vspace{1cm}
      \begin{itemize}
        \item 10-fold cross validation over a grid of bandwidth and type
    \item Boxcar, 136, has the lowest RMSE. Mean absolute error (MAE) gives similar results
      \item Bias variance trade off: test error is convex and training error decreases with bandwith
    \end{itemize}
    }
  \end{minipage}
  
\end{frame}

\begin{frame}{Out-of-sample predictions: actual versus predicted}
  \begin{figure}
    \includegraphics[trim={0 0 0 0},clip, scale = 0.5]{\figs/actual_predicted.png}
  \end{figure}
\end{frame}

\begin{frame}{Out of sample predictions: overall model performance}
\begin{minipage}[t]{0.58\linewidth}%    
  {\scriptsize
    % latex table generated in R 3.6.3 by xtable 1.8-3 package
    % Tue Mar  3 16:27:28 2020
    \begin{table}[ht]
      \centering
      \begin{tabular}{p{1.9cm}llll}
        \hline
        & & & & \\
         & RMSE & MAE & $\overline{\text{SRMSE}}$ & $\overline{\text{MASE}}$ \\
        \hline
        {OLS  \\ (bisquare, 340)} & 61.0427 & 14.0468 & 0.0304 & 0.4318 \\ 
        {Poisson \\ (boxcar, 136)} & 34.8364 & 11.7120 & 0.0260 & 0.4013 \\ 
        \hline
      \end{tabular}
    \end{table}
  }
  \flushleft
  \begin{align*}

    &\text{Where scaled errors are defined as:} \\
    &\text{} \\
      & SRMSE_{i} = \frac{RMSE_{i}}{\sum_{j} M_{i,j}} \\
      & MASE_{i} = \frac{1}{J} \sum\limits_{j=1}^{J} \lvert \frac{e_{i,j}}{\frac{1}{J} \sum_{j=1}^{J} \lvert M_{i,j} - \bar{M}_{i} \rvert} \rvert
  \end{align*}
\end{minipage}
\begin{minipage}[t]{0.38\linewidth}%
  \begin{figure}
    \includegraphics[trim={0 0 0 0},clip, scale = 0.45]{\figs/mase.png}
  \end{figure}
\end{minipage}
\end{frame}

\begin{frame}{Estimated parameters}
  \begin{minipage}[t]{0.58\linewidth}%
    \vspace{-0.5cm}
  \begin{figure}
    \includegraphics[trim={0 0 0 0},clip, scale = 0.4]{\figs/coeffs.png}
  \end{figure}
  \end{minipage}
  \begin{minipage}[t]{0.38\linewidth}%
    \vspace{1cm}
  \begin{figure}
    \includegraphics[trim={0 0 0 0},clip, scale = 0.5]{\figs/tvals.png}
  \end{figure}
 \end{minipage}
\end{frame}

\begin{frame}{Conclusions and further work}
  \begin{itemize}
    \item Estimated parameters exhibit spatial non-stationarity and cross validation suggests there is a bias variance trade off: need for a spatially explicit model
    \item Poisson estimator has smaller out of sample errors than the OLS estimator: reduces overprediction of large flows
      \item Fixed effects effectively ensure that origin constraints hold, also with GWR
      \item Future work 1: improvement of the dwelling variable
        \item Future work 2: is 35 km a good cut off for short- versus long-distance?
        \item Potential future work 3: evidence of overdispersion, information criteria suggest that Negative Binomial model is more appropriate. But predictions are terrible!
          \item Up for a challenge? Code and data online\footnote{\scriptsize\url{https://github.com/trondhusby/gwr}}
    \end{itemize}
  \end{frame}

%--- References --------------------------------%
\begin{frame}[t,allowframebreaks]{References}
\printbibliography
\end{frame}


\begin{frame}{Appendix: proof that Poisson ensures origin constraint \citep{davies1987statistical}}
  \scriptsize
  \begin{align*}
    &\text{Log-likelihood of Poisson (ignoring the constant) is given by} \\
    &\ell = \sum\limits_{i} \sum\limits_{j} [ M_{i,j} x^{T}_{i,j} \beta - \exp(x^{T}_{i,j}\beta) ] \\
    &\text{where $\beta$ is a vector of parameters. The derivative wrt to the $u^{th}$ structural parameter is} \\
    &\frac{\partial \ell}{\partial \beta_{u}} = \sum\limits_{i} \sum\limits_{j} x^{T}_{i,j, u} \left[ M_{i,j} - \exp(x^{T}_{i,j}\beta)\right] \\
    &\text{Let the $r^{th}$ variable be a dummy variable if $i = I$. Then } \\
    &\frac{\partial \ell}{\partial \beta_{r}} = \sum\limits_{j} \left[ M_{I,j} - \exp(x^{T}_{I,j} \beta)\right] \\
    &\text{Maximum likelihood implies that derivatives are zero. \n Let $\hat{\beta}$ be the maximum likelihod estimates of $\beta$, then we have}\\
    &\sum\limits_{j} M_{I,j} = \sum\limits_{j} \exp(x^{T}_{I,j} \hat{\beta}) = \hat{\lambda}_{I,j}
    \Rightarrow \sum\limits_{j} M_{i,j} = \sum\limits_{j} \hat{\lambda}_{i,j} \quad \forall i
  \end{align*}
\end{frame}


%--- frame 9 --------------------------------%
\begin{frame}{Appendix: Origin constraints with fixed effects in GWR}
  Example: $\mathbf{WX}$ for three regions
  \begin{table}
    \begin{tabular}{lll}
      i & j &  $X$  \\
      1 & 2 &  $x_{1,2}$  \\
      1 & 3 &  $x_{1,3}$  \\
      2 & 1 &  $w_{1,2}x_{2,1}$  \\
      2 & 3 &  $w_{1,2}x_{2,3}$  \\
      3 & 1 &  $w_{1,3}x_{3,1}$  \\
      3 & 2 &  $w_{1,3}x_{3,2}$  \\
    \end{tabular}
    \qquad
  %\end{table} \quad
  %\begin{table}
    \begin{tabular}{lll}
      i & j &  $X$  \\
      1 & 2 &  $w_{2,1}x_{1,2}$  \\
      1 & 3 &  $w_{2,1}x_{1,3}$  \\
      2 & 1 &  $x_{2,1}$  \\
      2 & 3 &  $x_{2,3}$  \\
      3 & 1 &  $w_{2,3}x_{3,1}$  \\
      3 & 2 &  $w_{2,3}x_{3,2}$  \\
    \end{tabular}
    \qquad
  %\end{table} \quad
  %\begin{table}
    \begin{tabular}{lll}
      i & j &  $X$  \\
      1 & 2 &  $w_{3,1}x_{1,2}$  \\
      1 & 3 &  $w_{3,1}x_{1,3}$  \\
      2 & 1 &  $w_{3,2}x_{2,1}$  \\
      2 & 3 &  $w_{3,2}x_{2,3}$  \\
      3 & 1 &  $x_{3,1}$  \\
      3 & 2 &  $x_{3,2}$  \\
      \end{tabular}
    \end{table}
\end{frame}

%--- frame 10 --------------------------------%
\begin{frame}{Origin constraints with fixed effects in GWR}
  Zooming in on $i=1$: $FE1$ becomes a non-weighted intercept, and it ensures that the origin constraint holds! However at a significant computational cost...
  \begin{table}
    \begin{tabular}{llllll}
      $i$ & $j$ &  $X$ & $FE1$ & $FE2$ & $FE3$ \\
      1 & 2 &  $x_{1,2}$ & 1 & 0 & 0 \\
      1 & 3 &  $x_{1,3}$ & 1 & 0 & 0 \\
      2 & 1 &  $w_{1,2}x_{2,1}$ & 0 & $w_{1,2}$ & 0 \\
      2 & 3 &  $w_{1,2}x_{2,3}$ & 0 & $w_{1,2}$ & 0 \\
      3 & 1 &  $w_{1,3}x_{3,1}$ & 0 & 0 & $w_{1,3}$ \\
      3 & 2 &  $w_{1,3}x_{3,2}$ & 0 & 0 & $w_{1,3}$ \\
      \end{tabular}
  \end{table}
  
\end{frame}

%--- frame 11 --------------------------------%
\begin{frame}{Origin constraints with fixed effects in GWR}
  ... but we can exploit \textbf{sparsity} if the weight matrix is full of zeros. This is the case with non-continuous kernel types where bandwith is smaller than global. Setting $w_{1,3} = 0$  
  \begin{table}
    \begin{tabular}{llllll}
      i & j &  $X$ & FE1 & FE2 & FE3 \\
      1 & 2 &  $x_{1,2}$ & 1 & 0 & \color{gray}{0} \\
      1 & 3 &  $x_{1,3}$ & 1 & 0 & \color{gray}{0} \\
      2 & 1 &  $w_{1,2}x_{2,1}$ & 0 & $w_{1,2}$ & \color{gray}{0} \\
      2 & 3 &  $w_{1,2}x_{2,3}$ & 0 & $w_{1,2}$ & \color{gray}{0} \\
      \color{gray}{3} & \color{gray}{1} &  \color{gray}{$w_{1,3}x_{3,1}$} & \color{gray}{0} & \color{gray}{0} & \color{gray}{$w_{1,3}$} \\
        \color{gray}{3} & \color{gray}{2} &  \color{gray}{$w_{1,3}x_{3,2}$} & \color{gray}{0} & \color{gray}{0} & \color{gray}{$w_{1,3}$} \\
      \end{tabular}
  \end{table}  
\end{frame}


\begin{frame}{Appendix: What drives the prediction errors? Road distance a potential candidate}
 \begin{minipage}[t]{0.48\linewidth}%
  \begin{figure}
    \includegraphics[trim={0 0 0 0},clip, scale = 0.5]{\figs/tvals.png}
  \end{figure}
  \footnotesize
  Distance is, in general, the most important variable...
 \end{minipage}
 \begin{minipage}[t]{0.48\linewidth}%
   \vspace{-0.5cm}
   \begin{figure}
    \includegraphics[trim={0 0.5cm 0 0},clip, scale = 0.4]{{/home/trond/Documents/PEARL/figs/gwr-outlier-dist-fig-1.png}}
   \end{figure}
   \footnotesize
   ...but road distance can be problematic. Lines in the figure are straight lines between two municipality centroids
 \end{minipage}
\end{frame}

\begin{frame}{Departures and arrivals}
  \begin{figure}
    \includegraphics[trim={0 0 0 0},clip, scale = 0.4]{\figs/vars.png}
  \end{figure}
\end{frame}

\begin{frame}{Explanatory variables, average value by destination}
  \begin{figure}
    \includegraphics[trim={0 0 0 0},clip, scale = 0.4]{\figs/arrivals_and_departures.png}
  \end{figure}
\end{frame}

\begin{frame}{Core-periphery patterns of distance decay and population}
  \begin{minipage}[t]{0.58\linewidth}%
    \begin{figure}
      \includegraphics[trim={0 0 0 0},clip, scale = 0.5]{\figs/distance_decay.png}
    \end{figure}
  \end{minipage}
  \begin{minipage}[t]{0.38\linewidth}%
    {\footnotesize
    \begin{itemize}
    \item Flows to centrally located areas are characterised by high sensitivity to distance and low sensitivity to population
    \item The opposite for flows within peripheral areas
    \item Similar results were obtained for Japan \citep{nakaya2001local}
    \end{itemize}
    }
  \end{minipage}
\end{frame}

%% \begin{frame}{Research Question (old)}
%%   %\small
%%   \begin{itemize}
%%   \item Earlier work focused on local modelling of origin-constrained gravity with \textit{geographically weighted regression} \citep{loke2012} 
%%     \item Empirical strategy: take logs of both sides, rewrite, and estimate with OLS \citep{fotheringham1989spatial}. Problematic because
%%       \begin{itemize}
%%       \item Ad-hoc solutions for zero flows \citep{burger2009specification}
%%       \item We are not per se interested in $\log \hat{M}_{ij}$, antilogarithms produce biased estimates of $\hat{M}_{ij}$ \citep{flowerdew1982method}
%%       \item In the presence of heteroskedasticity, estimated elasticities can be misleading \citep{silva2006log}
%%       \end{itemize}
%%   \item \textbf{Constrained gravity as a local count model: which specification minimises \textit{out-of-sample} prediction errors?}
%%   \end{itemize}
  
%% \end{frame}

\end{document}

